\documentclass[14pt,a4paper,oneside]{extbook}
\usepackage[utf8]{vietnam}
\usepackage{amsmath,amssymb,mathrsfs,fancyhdr,enumerate,multirow,makecell}
\usepackage{tkz-euclide,tikz-3dplot,tikz,tkz-tab}
\usetikzlibrary{arrows.meta,calc,intersections,angles,quotes,backgrounds,shapes.geometric,patterns,through,shadows,positioning}
\usepackage{yhmath}
\usepackage[top=0.7cm, bottom=0.7cm, left=1.0cm, right=1.0cm, headheight=16pt, headsep=3mm, footskip=5mm]{geometry}
\usepackage[hidelinks,unicode]{hyperref}
\usepackage{qrcode}
\usepackage[many]{tcolorbox}
\usepackage{fontawesome}

\pagestyle{fancy}
\fancyhf{}
\fancyhead[L]{\footnotesize\bfseries\color{teal!85!black}\faBook~SỔ TAY KIẾN THỨC TOÁN 4}
\fancyhead[R]{\footnotesize\bfseries\color{blue!85!black}Lớp Toán Thầy Lĩnh -- loptoanthaylinh.vn}
\fancyfoot[C]{\footnotesize\bfseries\color{gray!80!black}-- Sổ tay kiến thức: Bài 55. Tính chất cơ bản của phân số --}
\renewcommand{\headrulewidth}{0.4pt}

\begin{document}

% --- BANNER TIÊU ĐỀ BÀI HỌC ---
\noindent
\begin{tcolorbox}[
  enhanced,
  colback=teal!8!white,
  colframe=teal!80!black,
  arc=3mm, boxrule=1pt,
  left=3mm, right=3mm, top=2mm, bottom=2mm,
  halign=center
]
  {\Large\bfseries\color{teal!90!black}\faBookmark\ BÀI 55: TÍNH CHẤT CƠ BẢN CỦA PHÂN SỐ}\\[1.5pt]
  {\footnotesize\itshape\color{black!75} Quy tắc nhân chia cả tử và mẫu với cùng một số tự nhiên khác 0 trên 1 trang A4}
\end{tcolorbox}

\vspace{-2.5mm}

\noindent
\begin{minipage}[t]{0.49\textwidth}\vspace{0pt}
\begin{tcolorbox}[
  enhanced,
  colback=blue!4!white,
  colframe=blue!75!black,
  arc=2mm, boxrule=0.8pt,
  left=2mm, right=2mm, top=1.8mm, bottom=1.8mm,
  title={\footnotesize\bfseries\color{white}\faTimesCircle\ 1. Quy tắc Nhân (Nhân lên)}
]
{\footnotesize
Nếu \textbf{nhân cả tử số và mẫu số} của một phân số với cùng một số tự nhiên khác 0 thì được một \textbf{phân số bằng phân số đã cho}:
$$\dfrac{a}{b} = \dfrac{a \times m}{b \times m} \quad (m \ne 0)$$
\textbf{Ý nghĩa:}
\begin{itemize}\setlength{\itemsep}{1pt}\setlength{\leftskip}{-2.5mm}
  \item Chia nhỏ các phần hơn nhưng giá trị tổng thể không đổi.
  \item Là cơ sở trực tiếp của phương pháp \textbf{Quy đồng mẫu số}.
\end{itemize}
}
\end{tcolorbox}
\end{minipage}
\hfill
\begin{minipage}[t]{0.49\textwidth}\vspace{0pt}
\begin{tcolorbox}[
  enhanced,
  colback=teal!4!white,
  colframe=teal!75!black,
  arc=2mm, boxrule=0.8pt,
  left=2mm, right=2mm, top=1.8mm, bottom=1.8mm,
  title={\footnotesize\bfseries\color{white}\faMinusCircle\ 2. Quy tắc Chia (Chia bớt)}
]
{\footnotesize
Nếu \textbf{chia cả tử số và mẫu số} của một phân số cho cùng một số tự nhiên khác 0 (ước chung) thì được một \textbf{phân số bằng phân số đã cho}:
$$\dfrac{a}{b} = \dfrac{a : n}{b : n} \quad (n \ne 0)$$
\textbf{Ý nghĩa:}
\begin{itemize}\setlength{\itemsep}{1pt}\setlength{\leftskip}{-2.5mm}
  \item Gộp các phần lại giúp phân số gọn gàng hơn.
  \item Là cơ sở trực tiếp của phương pháp \textbf{Rút gọn phân số}.
\end{itemize}
}
\end{tcolorbox}
\end{minipage}

\vspace{-2.5mm}
% --- KHỐI 3: SƠ ĐỒ TƯ DUY TIKZ ---
\noindent
\begin{tcolorbox}[
  enhanced,
  colback=orange!4!white,
  colframe=orange!80!black,
  arc=2mm, boxrule=0.8pt,
  left=2mm, right=2mm, top=1.5mm, bottom=1.5mm,
  title={\footnotesize\bfseries\color{white}\faSitemap\ SƠ ĐỒ TƯ DUY: HAI CHIỀU TÍNH CHẤT CƠ BẢN CỦA PHÂN SỐ}
]
\centering
\begin{tikzpicture}[
  box/.style={draw, rounded corners=2pt, align=center, font=\scriptsize\bfseries, inner sep=3pt},
  arr/.style={-{Stealth[length=2.5mm]}, thick}
]
  \node[box, fill=purple!12, draw=purple!80!black, text width=3.6cm, minimum height=1.2cm] (mid) at (5.2, 0) {PHÂN SỐ GỐC\\$\mathbf{\dfrac{a}{b}}$ ($b \ne 0$)};

  \node[box, fill=teal!12, draw=teal!80!black, text width=4.0cm, minimum height=1.2cm] (left) at (0.2, 0) {RÚT GỌN PHÂN SỐ\\$\mathbf{\dfrac{a : n}{b : n}}$ ($n$ là ước chung)\\Tử \& Mẫu cùng chia $n$};

  \node[box, fill=blue!12, draw=blue!80!black, text width=4.0cm, minimum height=1.2cm] (right) at (10.2, 0) {QUY ĐỒNG PHÂN SỐ\\$\mathbf{\dfrac{a \times m}{b \times m}}$ ($m \ne 0$)\\Tử \& Mẫu cùng nhân $m$};

  \draw[arr, teal!80!black] (mid.west) -- (left.east) node[midway, above, font=\tiny\bfseries] {Chia cả tử \& mẫu};
  \draw[arr, blue!80!black] (mid.east) -- (right.west) node[midway, above, font=\tiny\bfseries] {Nhân cả tử \& mẫu};
\end{tikzpicture}
\end{tcolorbox}

\vspace{-2.5mm}
% --- KHỐI 4: VÍ DỤ & MẸO TRÁNH BẪY ---
\noindent
\begin{minipage}[t]{0.49\textwidth}\vspace{0pt}
\begin{tcolorbox}[
  enhanced,
  colback=green!4!white,
  colframe=green!60!black,
  arc=2mm, boxrule=0.8pt,
  left=2mm, right=2mm, top=1.8mm, bottom=1.8mm,
  title={\footnotesize\bfseries\color{white}\faCheckCircle\ Ví dụ thực hành chuẩn mực}
]
{\footnotesize
\textbf{1. Tìm phân số bằng nhau (Nhân lên):}\\
$\dfrac{2}{3} = \dfrac{2 \times 4}{3 \times 4} = \mathbf{\dfrac{8}{12}}; \quad \dfrac{3}{5} = \dfrac{3 \times 2}{5 \times 2} = \mathbf{\dfrac{6}{10}}$.\\
\textbf{2. Tìm phân số bằng nhau (Chia bớt):}\\
$\dfrac{15}{20} = \dfrac{15 : 5}{20 : 5} = \mathbf{\dfrac{3}{4}}; \quad \dfrac{18}{24} = \dfrac{18 : 6}{24 : 6} = \mathbf{\dfrac{3}{4}}$.
}
\end{tcolorbox}
\end{minipage}
\hfill
\begin{minipage}[t]{0.49\textwidth}\vspace{0pt}
\begin{tcolorbox}[
  enhanced,
  colback=red!4!white,
  colframe=red!70!black,
  arc=2mm, boxrule=0.8pt,
  left=2mm, right=2mm, top=1.8mm, bottom=1.8mm,
  title={\footnotesize\bfseries\color{white}\faExclamationTriangle\ Mẹo tránh bẫy trắc nghiệm}
]
{\footnotesize
\begin{itemize}\setlength{\itemsep}{1.5pt}\setlength{\leftskip}{-2.5mm}
  \item \textbf{Bẫy chỉ nhân/chia ở tử mà quên mẫu:}\\
  Phải nhân (hoặc chia) \textbf{cả tử số VÀ mẫu số} với cùng một số.
  \item \textbf{Bẫy cộng/trừ cùng một số:}\\
  Cộng hoặc trừ cùng một số vào tử và mẫu \textbf{KHÔNG} tạo ra phân số bằng nhau! (Ví dụ: $\frac{1+1}{2+1} = \frac{2}{3} \ne \frac{1}{2}$).
\end{itemize}
}
\end{tcolorbox}
\end{minipage}

\end{document}
